% §2 背景与动机
\label{sec:background}

\subsection{溯源图与APT检测}

溯源图 $\graph = (V, E)$ 是由内核级审计日志构建的有向图。每个节点 $v \in V$ 代表一个系统实体（进程、文件或网络套接字），每条有向边 $e = (u \rightarrow v, t, \tau) \in E$ 代表一个因果系统事件——实体 $u$ 在时间戳 $t$ 通过事件类型 $\tau \in \{\text{FORK}, \text{EXEC}, \text{READ}, \text{WRITE}, \text{SENDTO}, \text{RECVFROM}, \dots\}$ 影响了实体 $v$。

溯源图具备两个使其在入侵检测中格外强大的特性。第一，它们是\emph{完备的}：由于内核调解所有系统交互，系统实体之间的每个因果关系都被记录下来。第二，它们是\emph{防篡改的}：攻击者无法在未获得内核级访问权限的情况下修改内核级审计日志——而一旦获得内核级访问，系统已经彻底沦陷。

溯源数据的规模带来了根本性挑战。DARPA透明计算E3数据集（一个标准基准）包含268,242个节点和29,727,441条边，跨越五天。从这样的体量中提取有意义的攻击信号需要有原则的过滤。

\subsection{检测粒度鸿沟}

已有PIDS在逐步细化的粒度上运行，形成一个谱系：

\begin{itemize}
\item \textbf{图级} (Unicorn~\cite{unicorn2020})：将整个溯源图分类为良性或恶意。
\item \textbf{快照级} (PROGRAPHER~\cite{yang2023prographer})：将图划分为时序快照并对每个快照进行分类。
\item \textbf{窗口级} (KAIROS~\cite{cheng2024kairos})：在15分钟时间窗口内对事件评分，标记异常密度高的窗口。
\item \textbf{实体级} (MAGIC~\cite{jia2024magic}, ThreaTrace~\cite{wang2022threatrace})：对单个系统实体（进程、文件）评分。
\item \textbf{边级} (EdgeTrace~\cite{edgetrace2025})：对单条溯源边进行异常评分。
\end{itemize}

每个级别回答了一个逐步细化的问题，但没有人回答对调查最重要的那个问题：\textbf{“哪条因果相连的事件序列构成了攻击？”} 我们称之为\emph{链级}检测。

\subsection{已有链提取方法的局限性}

近期的工作已经开始将Transformer应用于溯源图。尽管这些方法代表了重要进展，但它们的链提取策略共享一个共同的局限性：不尊重因果边方向。

\textbf{固定窗口 (EagleEye~\cite{eagleeye2024})。} EagleEye将溯源图分割为固定大小的时间窗口，并将每个窗口内的事件线性化为Transformer分类的序列。该方法有两个缺陷。首先，时间$T$处的固定窗口可能包含恰好同时发生的、来自不相关进程的事件——Firefox页面加载和SSH登录可能落入同一窗口但毫无因果关系。其次，跨越窗口边界的因果相连事件被切断。

\textbf{随机游走 (Sentient~\cite{sentient2026})。} Sentient使用基于随机游走的场景分割来为Graph Transformer提取子图。随机游走在本质上是方向无关的：一次游走可能沿逆向遍历溯源边，将读取文件的过程与写入文件的过程连接起来，而不尊重这些事件发生的时间顺序。

\textbf{时序子图 (GET-AID~\cite{getaid2025})。} GET-AID基于事件时间戳的接近程度构造时序子图。虽然时间戳提供了一个弱信号，但时间接近并不意味因果连接：两个同时发生的事件可能属于完全独立的进程树。

\textbf{后处理重建 (SLOT~\cite{slot2025}, CAGE~\cite{cage2025})。} 这些方法首先检测异常节点或边，然后通过聚类或因果加权将它们连接成攻击路径。链是后处理的产出物，而非检测单元本身。关键的是，如果上游检测器漏掉单个关键节点，重建的链就不完整——且没有机制可以恢复它。

\textbf{基于规则的提取 (SPARSE~\cite{sparse2024}, Holmes~\cite{holmes2019})。} 这些方法使用专家定义的规则和威胁情报对路径进行评分和提取。它们无法检测其模式不在预定义规则集内的新型攻击。

\subsection{缺失的机制}

所有上述方法共享一个基本做法：首先决定\emph{如何分割}溯源图（按时间、按随机游走、按规则），然后依赖下游模型来验证分割是否有用。没有方法首先定义\emph{什么属性使一条链因果相干}，然后提取满足这些属性的链。

这是我们要填补的缺口。在\S\ref{sec:design}中，我们定义了一个量化链质量的因果相干性度量（独立于任何模型），并设计了一个最大化该度量的提取算法。
